\documentclass[12pt, letterpaper]{article}
\usepackage[margin=0.75in, top=0.75in, bottom=0.75in]{geometry}
\usepackage{booktabs}
\usepackage{longtable}
\usepackage{array}
\usepackage{xcolor}
\usepackage{colortbl}
\usepackage{hyperref}
\usepackage{helvet}
\renewcommand{\familydefault}{\sfdefault}

% Color definitions
\definecolor{headerbg}{RGB}{31, 73, 125}
\definecolor{headertext}{RGB}{255, 255, 255}
\definecolor{row1}{RGB}{235, 241, 250}
\definecolor{row2}{RGB}{255, 255, 255}
\definecolor{domainA}{RGB}{198, 217, 241}
\definecolor{domainB}{RGB}{217, 234, 211}
\definecolor{domainC}{RGB}{255, 230, 199}
\definecolor{domainD}{RGB}{234, 209, 220}

\hypersetup{
    colorlinks=true,
    linkcolor=blue,
    urlcolor=blue
}

\begin{document}

\begin{center}
    {\LARGE \textbf{3rd Grade Math --- Full Year Curriculum}}\\[4pt]
    {\large 25 Lessons across 11 Units}
\end{center}

\vspace{6pt}

\small
\begin{longtable}{
    >{\centering\arraybackslash}p{0.4cm}   % \#
    >{\raggedright\arraybackslash}p{3.0cm}  % Domain
    >{\raggedright\arraybackslash}p{2.8cm}  % Unit
    >{\raggedright\arraybackslash}p{3.2cm}  % Lesson Title
    >{\raggedright\arraybackslash}p{5.2cm}  % Description
    >{\centering\arraybackslash}p{2.0cm}    % Standard
}

% Header
\rowcolor{headerbg}
\textcolor{headertext}{\textbf{\#}} &
\textcolor{headertext}{\textbf{Domain}} &
\textcolor{headertext}{\textbf{Unit}} &
\textcolor{headertext}{\textbf{Lesson Title}} &
\textcolor{headertext}{\textbf{Description}} &
\textcolor{headertext}{\textbf{Standard}} \\
\endfirsthead

\rowcolor{headerbg}
\textcolor{headertext}{\textbf{\#}} &
\textcolor{headertext}{\textbf{Domain}} &
\textcolor{headertext}{\textbf{Unit}} &
\textcolor{headertext}{\textbf{Lesson Title}} &
\textcolor{headertext}{\textbf{Description}} &
\textcolor{headertext}{\textbf{Standard}} \\
\endhead

\multicolumn{6}{r}{\small\itshape Continued on next page\ldots} \\
\endfoot

\bottomrule
\endlastfoot

\toprule

% ── DOMAIN 1: Operations & Algebraic Thinking ──────────────────────────────────
\multicolumn{6}{l}{\cellcolor{domainA}\textbf{Domain: Operations \& Algebraic Thinking}} \\
\midrule

\rowcolor{row1}
1 & Operations \& Algebraic Thinking & Unit 1: Multiplication Concepts & Understanding Multiplication &
Understand multiplication as groups of equal objects, repeated addition, and grid arrays. Describe real-world multiplication contexts. &
3.OA.A.1 \\

\rowcolor{row2}
2 & Operations \& Algebraic Thinking & Unit 1: Multiplication Concepts & Multiplication on a Number Line &
Represent multiplication problems on a number line. Model jumps of equal sizes to calculate products. &
3.OA.A.3 \\

\rowcolor{row1}
3 & Operations \& Algebraic Thinking & Unit 1: Multiplication Concepts & Commutative \& Associative Properties &
Understand commutative ($a \times b = b \times a$) and associative properties of multiplication to multiply single-digit numbers. &
3.OA.B.5 \\

\rowcolor{row2}
4 & Operations \& Algebraic Thinking & Unit 2: Division Concepts & Understanding Division &
Understand division as sharing equally and partition grouping. Interpret quotients as shares or groups. &
3.OA.A.2 \\

\rowcolor{row1}
5 & Operations \& Algebraic Thinking & Unit 2: Division Concepts & Multiplication \& Division Relationships &
Explore the relationship between multiplication and division as inverse operations. Find the missing factor in division families. &
3.OA.B.6 \\

\rowcolor{row2}
6 & Operations \& Algebraic Thinking & Unit 3: Multiplication \& Division Facts & Facts for 2s, 5s, and 10s &
Develop speed and fluency with multiplication and division facts of 2s, 5s, and 10s using patterns and skip counting. &
3.OA.C.7 \\

\rowcolor{row1}
7 & Operations \& Algebraic Thinking & Unit 3: Multiplication \& Division Facts & Facts for 3s, 4s, and 6s &
Learn and practice multiplication and division facts for 3s, 4s, and 6s. Develop computational automaticity. &
3.OA.C.7 \\

\rowcolor{row2}
8 & Operations \& Algebraic Thinking & Unit 3: Multiplication \& Division Facts & Facts for 7s, 8s, and 9s &
Learn and practice multiplication and division facts for 7s, 8s, and 9s. Apply properties and tricks to aid facts recall. &
3.OA.C.7 \\

\rowcolor{row1}
9 & Operations \& Algebraic Thinking & Unit 4: Word Problems \& Patterns & Two-Step Word Problems &
Solve two-step word problems involving all four operations. Represent problems with equations and check solutions for reasonableness. &
3.OA.D.8 \\

\rowcolor{row2}
10 & Operations \& Algebraic Thinking & Unit 4: Word Problems \& Patterns & Identifying Arithmetic Patterns &
Identify arithmetic patterns in multiplication tables and addition charts. Use mathematical properties to explain patterns. &
3.OA.D.9 \\

\midrule
% ── DOMAIN 2: Number & Operations in Base Ten ──────────────────────────────────
\multicolumn{6}{l}{\cellcolor{domainB}\textbf{Domain: Number \& Operations in Base Ten}} \\
\midrule

\rowcolor{row1}
11 & Number \& Operations in Base Ten & Unit 5: Multi-Digit Arithmetic & Rounding Whole Numbers &
Round whole numbers to the nearest 10 or 100 using place value understanding. Apply rounding in real-world estimation. &
3.NBT.A.1 \\

\rowcolor{row2}
12 & Number \& Operations in Base Ten & Unit 5: Multi-Digit Arithmetic & Addition \& Subtraction within 1000 &
Fluently add and subtract whole numbers within 1000 using strategies based on place value, properties, and algorithms. &
3.NBT.A.2 \\

\rowcolor{row1}
13 & Number \& Operations in Base Ten & Unit 5: Multi-Digit Arithmetic & Multiplying by Multiples of 10 &
Multiply one-digit whole numbers by multiples of 10 (e.g., $9 \times 80$) using place value understanding and patterns. &
3.NBT.A.3 \\

\midrule
% ── DOMAIN 3: Number & Operations - Fractions ─────────────────────────────────
\multicolumn{6}{l}{\cellcolor{domainC}\textbf{Domain: Number \& Operations --- Fractions}} \\
\midrule

\rowcolor{row1}
14 & Number \& Operations --- Fractions & Unit 6: Introduction to Fractions & Fractions as Parts of a Whole &
Understand a fraction $1/b$ as the quantity formed by 1 part when a whole is partitioned into $b$ equal parts. Represent parts of a whole. &
3.NF.A.1 \\

\rowcolor{row2}
15 & Number \& Operations --- Fractions & Unit 6: Introduction to Fractions & Fractions on a Number Line &
Represent fractions on a number line diagram. Identify unit fraction intervals and plot numbers between 0 and 1. &
3.NF.A.2 \\

\rowcolor{row1}
16 & Number \& Operations --- Fractions & Unit 7: Fraction Equivalence & Equivalent Fractions &
Understand fraction equivalence and identify equivalent fractions. Represent whole numbers as fractions (e.g., $3 = 3/1$). &
3.NF.A.3a-c \\

\rowcolor{row2}
17 & Number \& Operations --- Fractions & Unit 7: Fraction Equivalence & Comparing Fractions &
Compare two fractions with the same numerator or same denominator using area models, number lines, and comparisons ($<$, $>$, $=$). &
3.NF.A.3d \\

\midrule
% ── DOMAIN 4: Measurement & Data ───────────────────────────────────────────────
\multicolumn{6}{l}{\cellcolor{domainD}\textbf{Domain: Measurement \& Data}} \\
\midrule

\rowcolor{row1}
18 & Measurement \& Data & Unit 8: Time, Volume \& Mass & Telling Time \& Elapsed Time &
Tell and write time to the nearest minute. Measure elapsed time intervals in minutes. Solve real-world word problems. &
3.MD.A.1 \\

\rowcolor{row2}
19 & Measurement \& Data & Unit 8: Time, Volume \& Mass & Measuring Liquid Volume \& Mass &
Measure and estimate liquid volumes (liters) and masses (grams, kilograms). Solve word problems using addition and subtraction. &
3.MD.A.2 \\

\rowcolor{row1}
20 & Measurement \& Data & Unit 9: Data \& Graphing & Picture Graphs \& Bar Graphs &
Draw scaled picture graphs and scaled bar graphs to represent a data set. Solve "how many more" and "how many less" problems. &
3.MD.B.3 \\

\rowcolor{row2}
21 & Measurement \& Data & Unit 9: Data \& Graphing & Measurement Data on Line Plots &
Generate measurement data by measuring lengths using rulers marked with halves and fourths of an inch. Show data on line plots. &
3.MD.B.4 \\

\rowcolor{row1}
22 & Measurement \& Data & Unit 10: Area \& Perimeter & Understanding Area \& Unit Squares &
Recognize area as an attribute of plane figures and understand concepts of area measurement. Measure area by counting unit squares. &
3.MD.C.5 / 3.MD.C.6 \\

\rowcolor{row2}
23 & Measurement \& Data & Unit 10: Area \& Perimeter & Area of Rectangles \& Decomposing &
Relate area to multiplication. Solve real-world area problems for rectangles, apply distributive property, and decompose irregular shapes. &
3.MD.C.7 \\

\rowcolor{row1}
24 & Measurement \& Data & Unit 10: Area \& Perimeter & Perimeter \& Area Relationships &
Solve mathematical problems involving perimeters of polygons. Find perimeters given side lengths, find missing sides, and compare areas. &
3.MD.D.8 \\

\midrule
% ── DOMAIN 5: Geometry & Financial Literacy ───────────────────────────────
\multicolumn{6}{l}{\cellcolor{domainA}\textbf{Domain: Geometry \& Financial Literacy}} \\
\midrule

\rowcolor{row2}
25 & Geometry \& Financial Literacy & Unit 11: Shape Classes \& Partitioning & Shape Classes \& Financial Literacy &
Classify shapes by attributes (e.g., quadrilaterals). Partition shapes into parts with equal areas. Understand money, coins, and bills. &
3.G.A.1 / 3.G.A.2 / 3.MD.D.8 (applied) \\

\bottomrule
\end{longtable}

\end{document}
